% Figure -- the Tabasco REPA setup (model construction + encoders tried).
% \input{} from Ch5 (sec:tabasco-models). Inline TikZ;
% needs \usetikzlibrary{arrows.meta,fit} (loaded in the preamble).
\begin{figure}[tbp]
\centering
\begin{tikzpicture}[
    >={Latex[length=2mm]},
    trunk/.style={rectangle, draw, rounded corners=2pt, minimum height=11mm, minimum width=32mm, align=center},
    enc/.style={rectangle, draw, rounded corners=2pt, fill=gray!18, minimum height=8mm, minimum width=36mm, align=center, font=\small},
    op/.style={rectangle, draw, rounded corners=2pt, minimum height=8mm, minimum width=16mm, align=center, font=\small},
    lbl/.style={font=\itshape, align=center},
    loss/.style={font=\bfseries, align=center},
    note/.style={font=\footnotesize, align=left},
    arr/.style={-{Latex[length=2mm]}, thick}
]
% top row: generation path
\node[lbl] (xin) at (0,0) {Noised\\molecule $x_t$};
\node[trunk] (trunk) at (3.8,0) {Tabasco trunk\\$\ell_1 \to \cdots \to \ell_L$};
\node[lbl] (vout) at (6.4,0) {velocity\\$v_\theta(x_t,t)$};
\node[op] (mse) at (8.4,0) {MSE};
\node[loss] (lfm) at (10.7,0) {$\mathcal{L}_{\text{FM}}$};
% projector tap (aligned column at x=3.8)
\node[op] (proj) at (3.8,-2.4) {Projector $h_\phi$};
\node[lbl] (pemb) at (3.8,-3.6) {$h_\phi(z_\ell)$};
% encoder stack (swappable alignment target)
\node[enc] (cm) at (3.8,-5.1) {CheMeleon \textit{(2D graph)}};
\node[enc] (mc) at (3.8,-6.2) {MACE \textit{(3D local)}};
\node[draw, dashed, rounded corners=3pt, fit=(cm)(mc), inner sep=2.5mm] (box) {};
% routing annotations
\node[note,anchor=west] at (5.95,-5.1) {$\rightarrow$ no transfer};
\node[note,anchor=west] at (5.95,-6.2) {$\rightarrow$ no transfer};
% clean molecule input
\node[lbl] (x1) at (0,-5.65) {Clean\\molecule $x_1$};
% cosine similarity (intermediate box)
\node[op] (sim) at (3.8,-7.7) {cosine sim};
% REPA loss terminal
\node[loss] (lrepa) at (7.3,-7.7) {$\mathcal{L}_{\text{REPA}}$};
% combined objective (mid-right)
\node[draw, rounded corners=2pt, inner sep=5pt, font=\small\bfseries, align=center]
  (ltotal) at (12.4,-3.85) {$\mathcal{L} = \mathcal{L}_{\text{FM}} + \lambda\,\mathcal{L}_{\text{REPA}}$};
% arrows
\draw[arr] (xin) -- (trunk);
\draw[arr] (trunk) -- (vout);
\draw[arr] (vout) -- (mse);
\draw[arr] (mse) -- (lfm);
\draw[arr] (trunk) -- (proj) node[midway,right]{\small $z_\ell$};
\draw[arr] (proj) -- (pemb);
\draw[arr] (pemb) -- (box.north);
\draw[arr] (x1) -- (x1-|box.west);
\draw[arr] (box.south) -- (sim);
\draw[arr] (sim) -- (lrepa);
\draw[arr] (lfm) -| (ltotal);
\draw[arr] (lrepa) -| (ltotal);
\end{tikzpicture}
\caption{The Tabasco REPA setup. The trunk is trained with flow matching ($\mathcal{L}_{\text{FM}}$, a mean-squared error on the velocity); its layer-$\ell$ hidden state is projected and aligned by cosine similarity to one frozen encoder per run, giving the REPA loss $\mathcal{L}_{\text{REPA}}$. The two combine into the training objective. Neither encoder transfers to generation.}
\label{fig:tabasco-setup}
\end{figure}
